% ============================================================
%  Briefing interno — Reunión con Dr. Fernando D. González-Nilo
%  Autor: Enzo Cevo
%  Uso: chuleta previa a la reunión (NO entregar al Dr. González-Nilo)
%  Compilar: pdflatex 2026-04-18_gonzalez-nilo-UNAB.tex  (x2 para refs)
% ============================================================
\documentclass[11pt,a4paper]{article}

\usepackage[utf8]{inputenc}
\usepackage[T1]{fontenc}
\usepackage[spanish,es-tabla]{babel}
\usepackage[margin=2cm]{geometry}
\usepackage{xcolor}
\usepackage{enumitem}
\usepackage{titlesec}
\usepackage{fancyhdr}
\usepackage{lastpage}
\usepackage{array}
\usepackage{tabularx}
\usepackage{booktabs}
\usepackage{hyperref}
\hypersetup{colorlinks=true,urlcolor=blue!60!black,linkcolor=black}

% Placeholders visibles en rojo
\newcommand{\ph}[1]{\textcolor{red}{\textbf{[#1]}}}
\newcommand{\pregunta}[1]{\item #1}
\newcommand{\prioritaria}[1]{\item \textbf{\textcolor{blue!70!black}{#1}}}

% Secciones compactas
\titlespacing*{\section}{0pt}{8pt}{4pt}
\titlespacing*{\subsection}{0pt}{6pt}{3pt}
\setlist{nosep,leftmargin=*,topsep=2pt}

% Header/footer
\pagestyle{fancy}
\fancyhf{}
\fancyhead[L]{\small Briefing interno — Reunión F.\ González-Nilo (UNAB)}
\fancyhead[R]{\small 2026-04-18}
\fancyfoot[L]{\small\textit{Uso interno — no entregar}}
\fancyfoot[R]{\small\thepage/\pageref{LastPage}}
\renewcommand{\headrulewidth}{0.4pt}

\begin{document}

\begin{center}
    {\Large \textbf{Reunión con Dr.\ Fernando D.\ González-Nilo}}\\[2pt]
    {\small Doctorado en Bioinformática y Biología de Sistemas --- UNAB / CBIB}\\[2pt]
    {\small Fecha: 2026-04-18 $\cdot$ Lugar: UNAB $\cdot$ Duración estimada: 60 min}
\end{center}

\vspace{-6pt}
\hrule
\vspace{4pt}

% ============================================================
\section*{\S1.\quad Pitch personal (60 segundos)}

\begin{itemize}
    \item Soy \textbf{Enzo Cevo}, \textbf{Licenciado en Biología} de la \textbf{Universidad de Chile}, especializado en \textbf{dinámica molecular clásica (MD all-atom)} aplicada a sistemas blandos y biomoleculares.
    \item Trabajo hoy en tres frentes complementarios: (i)~\textbf{membranas biomiméticas} (bicapas lipídicas y cristales líquidos liotrópicos con GROMACS/OPLS-aa), (ii)~\textbf{materiales fotónicos} (luminógenos AIE en matrices organizadas) y (iii)~\textbf{allostery en proteínas} (MD + MMPBSA con AMBER24 sobre una quinasa arqueal).
    \item Tengo \textbf{tres manuscritos en curso} donde soy responsable de la parte de MD y el análisis de trayectorias --- no son papers pendientes, están en escritura / revisión.
    \item Busco hacer el doctorado en bioinformática en UNAB porque quiero \textbf{escalar desde MD descriptiva hacia MD predictiva} (drug design, enhanced sampling, integración con ML) dentro de un grupo transdisciplinar con validación experimental --- exactamente lo que hace el CBIB.
    \item Lo que aporto: rodaje operativo en GROMACS y AMBER, análisis en Python (MDAnalysis, MMPBSA), experiencia colaborando con grupos experimentales (\textsuperscript{2}H-NMR, cristalografía X, espectroscopía de fluorescencia).
\end{itemize}

% ============================================================
\section*{\S2.\quad Tres manuscritos en curso}

\subsection*{2.1 --- AA Dynamics en bicapa SDS/DeOH (\textsuperscript{2}H-NMR + MD)}
\begin{itemize}
    \item \textbf{Título:} \textit{Location, Average Orientation and Dynamics of Amino Acids in a Membrane Mimetic: \textsuperscript{2}H-NMR and MD Study}.
    \item \textbf{Coautores:} A.\ San Martín, \textbf{E.\ Cevo}, B.\ Weiss-López, D.\ Muñoz-Gacitúa.
    \item \textbf{Sistema:} cristal líquido liotrópico discótico (DNLLC) SDS/DeOH/KCl, 274 SDS + 80 DeOH + 15\,543 aguas; aminoácidos deuterados ALA, GLU, LEU, MET.
    \item \textbf{Mi contribución:} montaje con Packmol, MD all-atom en \textbf{GROMACS 2021.1} (OPLS-aa, TIP4P, PME, LINCS), producción NPT de \textbf{200 ns por aminoácido}; análisis en Python con MDAnalysis (perfiles de densidad parcial, RDF, desplazamiento $\Delta Z$ como indicador de orientación).
    \item \textbf{Hallazgo punch:} todos los AA se anclan vía zwitterion; ALA/GLU quedan paralelos a la interfaz ($\Delta Z=0$); LEU/MET penetran $\sim$0{,}5~nm al core hidrofóbico, lo que explica atomísticamente la caída de splittings cuadrupolares C\textsubscript{1}--C\textsubscript{4} de SDS observada en \textsuperscript{2}H-NMR.
    \item \textbf{Estado:} en escritura, con análisis pendientes del perfil de densidad de MET y RDF final del thioether.
\end{itemize}

\subsection*{2.2 --- AIE Luminogens en cristal líquido TTAB/DeOH (P01)}
\begin{itemize}
    \item \textbf{Título:} \textit{Substituent-Controlled Penetration of AIE Luminogens into Liquid Crystal Media}.
    \item \textbf{Coautores:} F.\ Llana, F.\ Bravo, \textbf{E.\ Cevo}, M.\ Cecilia, I.\ Osorio-Román, B.\ Weiss-López, B.\ M.\ Squeo, C.\ Segura. Colaboración internacional con \textbf{CNR-SCITEC (Milán, Italia)} + UACh + UChile.
    \item \textbf{Sistema:} luminógenos AIE donor--acceptor con núcleo thiophene-aldehyde y sustituyentes variados (TTG, TTC, TTY, TTO) embebidos en cristal líquido liotrópico TTAB/DeOH.
    \item \textbf{Mi contribución:} MD all-atom de la partición del luminógeno en la bicapa, correlación entre el sustituyente (carbazol vs.\ trifenilamina, longitud de conjugación) y la profundidad de penetración / orientación; cruce con fotofísica (QY, lifetimes, Stokes shifts).
    \item \textbf{Hallazgo punch:} el substituyente controla si el luminógeno queda en la interfaz o se hunde en el core hidrofóbico, lo que determina su rendimiento de emisión AIE bajo confinamiento.
    \item \textbf{Estado:} escritura, versión v4.
\end{itemize}

\subsection*{2.3 --- MttPFK/GK: activación alostérica por AMP (bioinformática pura)}
\begin{itemize}
    \item \textbf{Título:} \textit{Structural Basis for AMP-Dependent Allosteric Activation of a Bifunctional ADP-Dependent PFK/GK from Methanothermococcus thermolithotrophicus}.
    \item \textbf{Colaboración:} grupo V.\ Guixé (UChile). Combinación de cristalografía X (sincrotrón MANACA, LNLS Brasil; 1{,}17--1{,}54~\AA) + ensayos cinéticos + MD + mutagénesis.
    \item \textbf{Mi contribución:} preparación de sistemas (Propka / PDB2PQR / tLEaP) para dos complejos (GK y PFK) con AMP alostérico; \textbf{MD en AMBER24 con ff19SB}, 5 réplicas $\times$ 400~ns por sistema; análisis \textbf{MMPBSA} (free energy binding global y por residuo); diseño in silico y simulación de mutantes Y56A, R58A.
    \item \textbf{Hallazgo punch:} AMP estabiliza una conformación \textbf{``super-cerrada''} catalíticamente competente nunca antes descrita en esta familia; los mutantes Y56A/R58A \emph{hiper-cierran} el dominio y atrapan sustrato (--40~kcal/mol vs.\ --15~WT para ADP), validando MMPBSA por residuo como herramienta predictiva para el diseño del knockout.
    \item \textbf{Por qué importa en esta conversación:} este paper es \textbf{la aplicación más bioinformática} de los tres --- allostery, MD de proteína, MMPBSA, archaea, diseño racional de mutantes. Es el puente directo con las líneas del CBIB.
    \item \textbf{Estado:} borrador con resultados completos, pendiente de discusión y figuras finales.
\end{itemize}

% ============================================================
\section*{\S3.\quad Perfil del Dr.\ González-Nilo (rápido)}

\begin{itemize}
    \item Director del \textbf{Centro de Bioinformática y Biología Integrativa (CBIB)}, Facultad de Ciencias de la Vida, UNAB.
    \item Director del \textbf{Doctorado en Bioinformática y Biología de Sistemas} (UNAB).
    \item Líneas: simulación molecular, bio-nanotecnología, canales iónicos, drug design, estructura y dinámica de proteínas.
    \item Enfoque metodológico del centro: ciclo iterativo \textbf{observación $\rightarrow$ modelado $\rightarrow$ simulación $\rightarrow$ validación experimental}, 9 PIs transdisciplinares.
    \item 90+ artículos ISI, 3 patentes.
    \item \textbf{Puntos de empalme directos} con mi trabajo: dinámica de proteínas (\S2.3), MD en membranas y materiales blandos (\S2.1, \S2.2), análisis de trayectorias en Python.
\end{itemize}

% ============================================================
\section*{\S4.\quad Preguntas para la reunión}
\textit{(preguntas prioritarias marcadas en \textbf{\textcolor{blue!70!black}{azul}})}

\subsection*{4.1 --- Sobre el programa de doctorado}
\begin{itemize}
    \prioritaria{¿Cuál es la \textbf{estructura curricular real} (cursos, candidatura, tesis) y la duración nominal vs.\ promedio efectivo de titulación?}
    \prioritaria{\textbf{Becas ANID Doctorado Nacional}: ¿cuántos cupos captura típicamente el programa? ¿Hay mecanismos internos UNAB en paralelo?}
    \pregunta{Requisitos formales de ingreso (GPA mínimo, publicaciones previas, examen, entrevista, idioma).}
    \pregunta{¿Cómo se asigna el tutor y cómo se compone el comité de tesis? ¿Se puede elegir tutor antes de postular?}
    \pregunta{¿Qué líneas tienen cupo disponible para la \textbf{cohorte 2026/2027}?}
    \pregunta{¿Existe modalidad \textbf{cotutela} internacional o con otras universidades chilenas (p.\ ej.\ UChile por mi colaboración previa)?}
    \pregunta{Dedicación: ¿full-time obligatoria, o compatible con TA / docencia parcial?}
\end{itemize}

\subsection*{4.2 --- Sobre sus proyectos en curso}
\begin{itemize}
    \prioritaria{¿Qué \textbf{Fondecyt / Anillo / Milenio} tiene vigente hoy y cuál es el problema concreto que están atacando (canales iónicos, drug design, otro)?}
    \prioritaria{¿Hay proyectos en \textbf{bio-nanotecnología con matrices lipídicas o poliméricas} donde mi experiencia en luminógenos AIE en cristal líquido / PMMA empalme directo?}
    \pregunta{\textbf{Colaboraciones wet activas}: ¿qué laboratorios validan experimentalmente sus predicciones? (muy importante para mí, dado que los 3 papers son colaboraciones MD$\leftrightarrow$experimento).}
    \pregunta{\textbf{Infraestructura HPC del CBIB}: colas, software licenciado (GROMACS, AMBER, NAMD, Schrödinger), acceso para tesistas.}
    \pregunta{Herramienta / forcefield de referencia del grupo (CHARMM-GUI, AMBER/ff19SB, GROMACS/CHARMM36m, Martini, OPLS-aa).}
    \pregunta{\textbf{ML / IA} en el grupo: ¿se usa para análisis de trayectorias, screening o diseño de ligandos? ¿Hay espacio para desarrollar esa línea?}
    \pregunta{¿Usan \textbf{enhanced sampling} (metadynamics, umbrella sampling, REMD) o \textbf{QM/MM} de forma rutinaria?}
\end{itemize}

\subsection*{4.3 --- Sobre el encaje y mi contribución}
\begin{itemize}
    \prioritaria{Dado mi expertise en MD de \textbf{luminógenos AIE en medios organizados} (\S2.2) y en \textbf{allostery con MMPBSA} (\S2.3), ¿ve una línea natural de continuidad hacia \textbf{drug-delivery, nanovehículos fluorescentes o biosensores} en su grupo?}
    \prioritaria{¿Qué \textbf{skill concreta} le falta hoy al grupo que yo debiera cerrar antes de postular? (candidatos que se me ocurren: MD dirigida de canales, enhanced sampling, QM/MM, ML sobre trayectorias).}
    \pregunta{¿Hay posibilidad de una \textbf{rotación / pasantía corta pre-postulación} (ej.\ 1--3 meses) para testear encaje antes de comprometernos mutuamente?}
    \pregunta{Expectativa razonable de productividad durante el doctorado: número de papers, proporción primer autor vs.\ colaboración.}
    \pregunta{¿Le interesa que le envíe los tres manuscritos en PDF tras la reunión para una lectura más a fondo?}
\end{itemize}

% ============================================================
\section*{\S5.\quad Matriz ``Qué puedo aportar''}

\renewcommand{\arraystretch}{1.25}
\begin{tabularx}{\textwidth}{@{}p{0.45\textwidth}X@{}}
\toprule
\textbf{Mi expertise actual} & \textbf{Línea del CBIB con la que empalma} \\
\midrule
MD all-atom en GROMACS (OPLS-aa, TIP4P) de lípidos y cristal líquido liotrópico. & Bio-nanotecnología, simulación de membranas y materiales blandos. \\
MD all-atom en AMBER24 (ff19SB) de proteína + cofactores, 5 réplicas $\times$ 400~ns. & Dinámica de proteínas, drug design, canales iónicos. \\
Análisis MMPBSA global y por residuo; diseño in silico de mutantes (Y56A, R58A) validado con activación experimental. & Diseño racional de ligandos, allostery, validación in silico previo a wet-lab. \\
Análisis de trayectorias en Python: MDAnalysis, RDF, perfiles de densidad, orden orientacional, $\Delta Z$, RMSD por réplica. & Pipelines reutilizables de análisis para los sistemas del centro. \\
Colaboración demostrada con grupos experimentales (\textsuperscript{2}H-NMR, cristalografía X sincrotrón, fotofísica de fluorescencia, cinética enzimática). & Encaje directo con el ciclo observación$\rightarrow$modelado$\rightarrow$simulación$\rightarrow$validación que define al CBIB. \\
Fondecyt Iniciación P02 activo (concentrador solar en PMMA, luminógenos AIE). & Red paralela de proyectos que puede derivar en colaboraciones cruzadas. \\
\bottomrule
\end{tabularx}

% ============================================================
\section*{\S6.\quad Cierre y próximos pasos}

\begin{itemize}
    \item[$\square$] Confirmar \textbf{fechas límite} de postulación (ANID Doctorado Nacional 2026/2027 y proceso interno UNAB).
    \item[$\square$] Pedir contacto de \textbf{1--2 estudiantes actuales} del doctorado para ``second opinion''.
    \item[$\square$] Acordar próximo paso concreto: visita al CBIB / charla interna / envío de CV + los tres manuscritos en PDF.
    \item[$\square$] \textbf{Email de seguimiento en $\leq$48~h} con resumen de lo acordado y deliverables pendientes de mi lado.
    \item[$\square$] Internamente: actualizar esta nota tras la reunión con las respuestas obtenidas.
\end{itemize}

\vfill
\begin{center}
    \small\textit{Documento interno de preparación --- Enzo Cevo --- \today}
\end{center}

\end{document}
